\uuid{ChUD}
\titre{ Forward pass d'un MLP : identification des erreurs }
\niveau{L3}
\module{Analyse de données}
\chapitre{Réseaux de neurones}
\sousChapitre{Forward pass, fonctions d'activation}
\theme{Forward pass, ReLU, régression, classification}
\auteur{Maxime Nguyen}
\datecreate{2026-03-19}
\organisation{AMSCC}
\difficulte{2}

\contenu{
\texte{
  On compare trois implémentations du forward pass d'un MLP à une couche cachée :
  \begin{itemize}
    \item[A.] \texttt{Z1 = X @ W1 + b1 ; A1 = relu(Z1) ; y\_hat = A1 @ W2 + b2}
    \item[B.] \texttt{A1 = relu(X) ; Z1 = A1 @ W1 + b1 ; y\_hat = Z1 @ W2 + b2}
    \item[C.] \texttt{Z1 = X @ W1 + b1 ; A1 = relu(Z1) ; y\_hat = relu(A1 @ W2 + b2)}
  \end{itemize}
}
\begin{enumerate}

\item
  \question{Lequel est correct pour une tâche de régression ?}
  \reponse{\textbf{A} est correct : pré-activation linéaire, activation ReLU, puis sortie linéaire (pas d'activation finale pour la régression).}

\item
  \question{Qu'est-ce qui ne va pas dans B ?}
  \reponse{
    B applique ReLU directement sur les entrées $X$ avant la première couche linéaire,
    ce qui n'a aucun sens : la pré-activation $Z^{[1]}$ doit être calculée \emph{avant} d'appliquer l'activation.
    De plus, $Z1 = A1 @ W1 + b1$ applique ensuite une deuxième transformation linéaire sans activation,
    rendant les deux premières couches entièrement linéaires.
  }

\item
  \question{C pourrait convenir pour quel type de problème ?}
  \reponse{
    C applique une ReLU en sortie, ce qui contraint $\hat{y} \geq 0$.
    Cela pourrait convenir pour une régression sur une cible strictement positive
    (prix, durée, quantité), mais rarement en pratique car cela limite l'expressivité du modèle.
    Pour la classification binaire, on utiliserait une sigmoïde, pas une ReLU.
  }

\end{enumerate}
}
